\section{Proposed Model}
\label{sec:model}

ECR-GNN is formulated as a compact family of reusable graph classification architectures. All variants share the same input projection, stacked message-passing layers, global mean pooling, and linear graph classifier; they differ only in how intermediate node states and graph-level summaries are reused across depth and across branches.

% Placeholder: insert a cleaned architecture overview figure here if a final diagram is prepared later.

\subsection{Operator Factory and Shared Backbone}

Let $\mathcal{G}=(\mathcal{V},\mathcal{E})$ be a graph with node feature matrix $\mathbf{X}\in\mathbb{R}^{n\times d_{\mathrm{in}}}$ and adjacency structure $\mathbf{A}$. The operator factory instantiates one of three message-passing layers,
\[
\Phi \in \{\mathrm{GCNConv},\ \mathrm{GATConv},\ \mathrm{TransformerConv}\},
\]
and reuses the same operator type throughout a model instance. For a hidden width $d$, every branch starts with an input projection
\begin{equation}
\mathbf{H}^{(0)} = \mathrm{ReLU}\!\left(\Phi_{\mathrm{in}}(\mathbf{X}, \mathbf{A})\right),
\end{equation}
followed by $L$ hidden layers and a readout
\begin{equation}
\mathbf{g} = \mathrm{MeanPool}\!\left(\mathbf{H}^{(L)}\right).
\end{equation}
The final graph prediction is produced by
\begin{equation}
\hat{\mathbf{y}} = \mathbf{W}_{\mathrm{clf}}\,\mathrm{Dropout}(\mathbf{g}) + \mathbf{b}_{\mathrm{clf}}.
\end{equation}

All six architectures differ only in how they update $\mathbf{H}^{(\ell)}$ and how they reuse graph-level embeddings across blocks.

\subsection{Single-Branch and Node-Level Residual Blocks}

\textbf{PlainGNN.} The plain baseline stacks $L$ copies of the same operator without residual reuse:
\begin{equation}
\mathbf{H}^{(\ell+1)} = \mathrm{Dropout}\!\left(\mathrm{ReLU}\!\left(\Phi_{\ell}\!\left(\mathbf{H}^{(\ell)}, \mathbf{A}\right)\right)\right).
\end{equation}
This model serves as the reference point for judging whether residual information flow is beneficial.

\textbf{NodeResGNN.} The first residual variant stores the previous hidden state within the same branch:
\begin{equation}
\mathbf{H}^{(\ell+1)} = \mathrm{Dropout}\!\left(\mathrm{ReLU}\!\left(\Phi_{\ell}\!\left(\mathbf{H}^{(\ell)} + \mathbf{H}^{(\ell-1)}, \mathbf{A}\right)\right)\right).
\end{equation}
This mechanism is lightweight and preserves historical features without creating extra branches.

\textbf{NodeCrossGNN.} The node-level cross-residual model maintains two parallel branches with separate input projections:
\begin{align}
\mathbf{H}^{(\ell+1)}_{1} &= \mathrm{Dropout}\!\left(\mathrm{ReLU}\!\left(\Phi^{(1)}_{\ell}\!\left(\mathbf{H}^{(\ell)}_{1} + \widetilde{\mathbf{H}}^{(\ell)}_{2}, \mathbf{A}\right)\right)\right), \\
\mathbf{H}^{(\ell+1)}_{2} &= \mathrm{Dropout}\!\left(\mathrm{ReLU}\!\left(\Phi^{(2)}_{\ell}\!\left(\mathbf{H}^{(\ell)}_{2} + \widetilde{\mathbf{H}}^{(\ell)}_{1}, \mathbf{A}\right)\right)\right),
\end{align}
where $\widetilde{\mathbf{H}}^{(\ell)}_{k}$ is the cached hidden state from the opposite branch before its current update. The final node representation is the branch sum
\begin{equation}
\mathbf{H}^{(L)} = \mathbf{H}^{(L)}_{1} + \mathbf{H}^{(L)}_{2}.
\end{equation}
Compared with the single-branch residual block, NodeCrossGNN explicitly exchanges information between parallel pathways at every depth.

\subsection{Graph-Level Residual Propagation}

We further consider three higher-level architectures that compose several block models sequentially and pass graph embeddings across them.

\textbf{GraphCondGNN.} Two plain message-passing blocks are executed in sequence. The first block produces a graph embedding $\mathbf{g}^{(1)}$, which is added to node features and to the pooled graph representation inside the second block:
\begin{equation}
\mathbf{g}^{(2)} = B_{2}\!\left(\mathbf{X}, \mathbf{A}, \mathbf{g}^{(1)}\right).
\end{equation}
This introduces graph-level conditioning without intra-layer residual connections.

\textbf{GraphResGNN.} This variant chains three graph-conditioned blocks:
\begin{equation}
\mathbf{g}^{(t+1)} = B_{t+1}\!\left(\mathbf{X}, \mathbf{A}, \mathbf{g}^{(t)}\right), \quad t \in \{1,2\}.
\end{equation}
Inside each block, the graph embedding is broadcast back to node states before the next graph convolution. The design can be viewed as repeated graph-level residual refinement.

\textbf{GraphCrossGNN.} The strongest graph-level interaction appears in the cross-graph variant. Four blocks are organized as two sequential pairs. Within each stage, the two block outputs are swapped and reused as the next stage's graph-level input:
\begin{align}
(\mathbf{g}^{(t)}_{1}, \mathbf{g}^{(t)}_{2}) &= \left(B_{2t-1}\!\left(\mathbf{X}, \mathbf{A}, \mathbf{g}^{(t-1)}_{1}\right),\ B_{2t}\!\left(\mathbf{X}, \mathbf{A}, \mathbf{g}^{(t-1)}_{2}\right)\right), \\
(\mathbf{g}^{(t+1)}_{1}, \mathbf{g}^{(t+1)}_{2}) &= \left(\mathbf{g}^{(t)}_{2},\ \mathbf{g}^{(t)}_{1}\right).
\end{align}
The final graph embedding is
\begin{equation}
\mathbf{g}_{\mathrm{final}} = \mathbf{g}^{(T)}_{1} + \mathbf{g}^{(T)}_{2}.
\end{equation}
In this variant, interaction occurs through exchanged pooled graph summaries rather than through direct node-tensor exchange.

\subsection{Unified Training Objective}

All model variants use the same graph classification objective:
\begin{equation}
\mathcal{L}(\theta) = - \frac{1}{N} \sum_{i=1}^{N}\sum_{c=1}^{C} y_{i,c}\log \hat{y}_{i,c},
\end{equation}
optimized with Adam. Dropout is applied after every message-passing update and once more before the classifier. The experiments in Section~\ref{sec:experiments} therefore isolate architectural differences rather than changes in loss, optimizer, or data protocol.

\subsection{Algorithmic View}

Algorithm~\ref{alg:ecr_forward} summarizes the shared forward logic in pseudocode form rather than as a descriptive table.

\begin{figure}[t]
\centering
\begin{minipage}{0.95\textwidth}
\begin{algorithmic}[1]
\STATE \textbf{Input:} graph batch $(\mathbf{X}, \mathbf{A}, \mathbf{b})$, operator $\Phi$, depth $L$, architecture type $M$
\STATE Instantiate message-passing blocks with operator $\Phi$
\IF{$M = \textsc{PlainGNN}$}
    \STATE Apply one input projection and stack $L$ message-passing layers
\ELSIF{$M = \textsc{NodeResGNN}$}
    \STATE Cache the previous hidden state and add it before each layer update
\ELSIF{$M = \textsc{NodeCrossGNN}$}
    \STATE Build two branches and inject the opposite branch's cached state at every depth
\ELSIF{$M \in \{\textsc{GraphCondGNN}, \textsc{GraphResGNN}\}$}
    \STATE Execute sequential block models and reuse the previous graph embedding as graph-level context
\ELSIF{$M = \textsc{GraphCrossGNN}$}
    \STATE Process block pairs, compute two graph embeddings, and swap them before the next pair
\ENDIF
\STATE Pool the final node representations with global mean pooling to obtain $\mathbf{g}$
\STATE Apply dropout and a linear classifier to produce logits $\hat{\mathbf{y}}$
\STATE \textbf{Return:} $(\hat{\mathbf{y}}, \mathbf{g})$
\end{algorithmic}
\end{minipage}
\caption{Pseudocode of the shared ECR-GNN forward procedure}
\label{alg:ecr_forward}
\end{figure}

\subsection{Design Implications}

This architecture family suggests three practical hypotheses that are later examined experimentally. First, residual reuse should matter more consistently than simply adding more branches, because every successful variant still relies on repeated feature reuse. Second, node-level cross residual and graph-level residual passing are not interchangeable: the former exchanges intermediate node representations, whereas the latter only re-injects pooled graph summaries. Third, the usefulness of cross-branch fusion depends on dataset scale and operator bias; it is therefore necessary to analyze accuracy, stability, and efficiency jointly rather than reporting only the single best score.
